% Options for packages loaded elsewhere
\PassOptionsToPackage{unicode}{hyperref}
\PassOptionsToPackage{hyphens}{url}
\documentclass[
  ignorenonframetext,
]{beamer}
\newif\ifbibliography
\usepackage{pgfpages}
\setbeamertemplate{caption}[numbered]
\setbeamertemplate{caption label separator}{: }
\setbeamercolor{caption name}{fg=normal text.fg}
\beamertemplatenavigationsymbolsempty
% remove section numbering
\setbeamertemplate{part page}{
  \centering
  \begin{beamercolorbox}[sep=16pt,center]{part title}
    \usebeamerfont{part title}\insertpart\par
  \end{beamercolorbox}
}
\setbeamertemplate{section page}{
  \centering
  \begin{beamercolorbox}[sep=12pt,center]{section title}
    \usebeamerfont{section title}\insertsection\par
  \end{beamercolorbox}
}
\setbeamertemplate{subsection page}{
  \centering
  \begin{beamercolorbox}[sep=8pt,center]{subsection title}
    \usebeamerfont{subsection title}\insertsubsection\par
  \end{beamercolorbox}
}
% Prevent slide breaks in the middle of a paragraph
\widowpenalties 1 10000
\raggedbottom
\AtBeginPart{
  \frame{\partpage}
}
\AtBeginSection{
  \ifbibliography
  \else
    \frame{\sectionpage}
  \fi
}
\AtBeginSubsection{
  \frame{\subsectionpage}
}
\usepackage{iftex}
\ifPDFTeX
  \usepackage[T1]{fontenc}
  \usepackage[utf8]{inputenc}
  \usepackage{textcomp} % provide euro and other symbols
\else % if luatex or xetex
  \usepackage{unicode-math} % this also loads fontspec
  \defaultfontfeatures{Scale=MatchLowercase}
  \defaultfontfeatures[\rmfamily]{Ligatures=TeX,Scale=1}
\fi
\usepackage{lmodern}
\usetheme[]{white}
\ifPDFTeX\else
  % xetex/luatex font selection
\fi
% Use upquote if available, for straight quotes in verbatim environments
\IfFileExists{upquote.sty}{\usepackage{upquote}}{}
\IfFileExists{microtype.sty}{% use microtype if available
  \usepackage[]{microtype}
  \UseMicrotypeSet[protrusion]{basicmath} % disable protrusion for tt fonts
}{}
\makeatletter
\@ifundefined{KOMAClassName}{% if non-KOMA class
  \IfFileExists{parskip.sty}{%
    \usepackage{parskip}
  }{% else
    \setlength{\parindent}{0pt}
    \setlength{\parskip}{6pt plus 2pt minus 1pt}}
}{% if KOMA class
  \KOMAoptions{parskip=half}}
\makeatother
\usepackage{longtable,booktabs,array}
\usepackage{caption}
\captionsetup[table]{skip=6pt}
\newcounter{none} % for unnumbered tables
\usepackage{calc} % for calculating minipage widths
\usepackage{caption}
% Make caption package work with longtable
\makeatletter
\def\fnum@table{\tablename~\thetable}
\makeatother
\setlength{\emergencystretch}{3em} % prevent overfull lines
\providecommand{\tightlist}{%
  \setlength{\itemsep}{0pt}\setlength{\parskip}{0pt}}
\usepackage{bookmark}
\IfFileExists{xurl.sty}{\usepackage{xurl}}{} % add URL line breaks if available
\urlstyle{same}
% fallback for those not using the hyperref driver hyperxmp:
\makeatletter
\@ifundefined{xmpquote}{\newcommand{\xmpquote}[1]{#1}}{}
\makeatother
\hypersetup{
  pdftitle={Help Desk Ticket SLA and Volume Forecasting},
  pdfauthor={SLA \& Volume Forecast Capstone},
  hidelinks,
  pdfcreator={LaTeX via pandoc}}

\title{\texorpdfstring{Help Desk Ticket SLA and Volume
Forecasting}{Help Desk Ticket SLA and Volume Forecasting}}
\subtitle{\texorpdfstring{Technical presentation --- peer/reviewer
deck}{Technical presentation --- peer/reviewer deck}}
\author{\texorpdfstring{SLA \& Volume Forecast
Capstone}{SLA \& Volume Forecast Capstone}}
\date{Historical analysis, 2016--2023}

\begin{document}
\frame{\titlepage}

\begin{frame}[fragile]{Help Desk Ticket SLA and Volume Forecasting}
\protect\phantomsection\label{help-desk-ticket-sla-and-volume-forecasting}
Technical peer deck --- reveal.js / Beamer-style export

\textbf{Source of truth:} this deck is a condensed, slide-formatted
derivative of \texttt{presentations/technical\_presentation.md}. That
narrative file (and the existing \texttt{technical\_capstone.pptx} deck)
are unchanged by this export.

Historical analysis of 27,348 scoped Ticket records, 2016-01-03 to
2023-03-14 UTC. Evidence boundary: the source ends in March 2023 and
does not validate current operations.
\end{frame}

\begin{frame}[fragile]
\begin{block}{1. Problem statement}
\protect\phantomsection\label{1-problem-statement}
\begin{itemize}
\tightlist
\item
  \textbf{Retrospective task:} compare elapsed resolution time of
  completed Ticket records with project-configured priority references.
\item
  \textbf{Predictive task:} forecast daily counts of exact
  \texttt{Ticket}-type issues, separate 7- and 30-day horizons, using
  only past counts and calendar features.
\item
  \textbf{ML framing:} univariate time-series regression via model
  comparison --- seasonal-naive baseline, weekly ETS, XGBoost
  autoregression (28-day lags), MI/PCA-selected Ridge. Supervised
  regression, not classification.
\item
  \textbf{Research acceptance criterion:} held-out MAE below the
  same-weekday baseline at both horizons --- met (5.05 vs 6.41 at 7
  days; 4.67 vs 6.49 at 30 days).
\item
  \textbf{Proposed business criterion (not measured/approved):} ≥10\%
  fewer overtime hours per 100 tickets in a human-reviewed pilot,
  without worsening an owner-approved service-attainment rate.
\end{itemize}
\end{block}
\end{frame}

\begin{frame}[fragile]
\begin{block}{2. Data provenance, licensing \& structure}
\protect\phantomsection\label{2-data-provenance-licensing--structure}
\begin{itemize}
\tightlist
\item
  Source: Help Desk Tickets (Mendeley Data v2, CC BY 4.0), mirrored on
  Kaggle under a conflicting CC0 label --- original attribution
  retained, mirror-label conflict flagged for resolution.
\item
  \texttt{data/raw/issues.csv}: 66,691 rows × 58 fields. Local file
  bytes verified against both the Kaggle mirror and the official
  Mendeley digest.
\item
  Documented scope is 2016--March 2023, but \texttt{issue\_created}
  contains dates back to 2007-04-15; pre-2016 rows are profiled and
  excluded from the primary analysis as a disclosed provenance caveat.
\item
  Primary cohort: exact \texttt{issue\_type\ ==\ Ticket}, 2016+ →
  \textbf{27,348 rows}; 344 of 2,628 calendar days (13.1\%) have no
  Ticket rows, filled as zero under an unverified completeness
  assumption.
\item
  Two structures are kept separate: the wide 58-column raw table
  (identity, actors, classification, timestamps, workflow-state fields)
  versus a narrow, purpose-built one-row-per-day forecasting table
  (target, past-only lags, rolling means, cyclical calendar encodings)
  --- this separation is what prevents same-day information leakage into
  the forecast features.
\end{itemize}
\end{block}
\end{frame}

\begin{frame}[fragile]
\begin{block}{3. Pipeline workflow}
\protect\phantomsection\label{3-pipeline-workflow}
\begin{itemize}
\tightlist
\item
  Eight notebooks run in sequence, each writing to one canonical
  \texttt{output/issue\_helpdesk/} folder: data understanding →
  EDA/feature engineering → model training/comparison → explainability
  (SHAP/PDP/ICE/ LIME) → bias/fairness audit → GenAI-scope dashboard
  contract → MLflow deployment/monitoring → business ROI template.
\item
  Downstream deliverables (decks, final report, FastAPI demo) all read
  from that same canonical output folder --- one reproducible source of
  truth.
\item
  Dashboard proof of work: a Power BI-ready data contract and folder
  (fact tables, model-comparison table, sample report screenshots) built
  without Power BI Desktop installed locally, substituted with rendered
  dashboard-styled proof images.
\item
  Full function/file-level mapping: \texttt{docs/process\_flow.md}.
\end{itemize}
\end{block}
\end{frame}

\begin{frame}
\begin{block}{4. Retrospective priority-SLA reference results}
\protect\phantomsection\label{4-retrospective-priority-sla-reference-results}
{\def\LTcaptype{none} % do not increment counter
\begin{longtable}[]{@{}lrrrr@{}}
\toprule\noalign{}
Mapped priority & Reference & Assessed & Met & Met rate \\
\midrule\noalign{}
\endhead
Critical & 4h & 2,124 & 524 & 24.7\% \\
High & 8h & 3,146 & 576 & 18.3\% \\
Medium & 24h & 11,161 & 2,161 & 19.4\% \\
Low & 24h & 304 & 72 & 23.7\% \\
\bottomrule\noalign{}
\end{longtable}
}

\begin{itemize}
\tightlist
\item
  16,735 eligible completed Tickets overall; \textbf{3,333 (19.9\%) met
  reference} --- consistent at roughly 1-in-5 across every priority,
  including Critical.
\item
  Duration is elapsed UTC wall-clock time, not confirmed business-hours
  or pause-adjusted SLA compliance --- a service owner must validate the
  clock and mapping before any compliance claim.
\end{itemize}
\end{block}
\end{frame}

\begin{frame}[fragile]
\begin{block}{5. From issue events to daily volume}
\protect\phantomsection\label{5-from-issue-events-to-daily-volume}
\begin{enumerate}
\tightlist
\item
  Parse \texttt{issue\_created} (UTC, mixed fractional-second formats).
\item
  Restrict to \texttt{issue\_type\ ==\ Ticket}, documented 2016+ scope.
\item
  Aggregate to one count per calendar day.
\item
  Fill no-row dates as zero, only under the stated completeness
  assumption.
\item
  Engineer past-only lag/rolling-mean/cyclical-calendar features.
\end{enumerate}

\begin{itemize}
\tightlist
\item
  Result: \textbf{2,628 daily observations}, averaging 10.4 tickets/day,
  344 assumed zero-arrival days, counts ranging 0--88.
\end{itemize}
\end{block}
\end{frame}

\begin{frame}
\begin{block}{6. Forecast experiment design}
\protect\phantomsection\label{6-forecast-experiment-design}
\begin{itemize}
\tightlist
\item
  7- and 30-day horizons, weekly-spaced \textbf{expanding-window}
  origins.
\item
  365-day validation period for model selection; separate 365-day final
  test period never used for selection.
\item
  7-day: 52 validation origins, 51 test origins. 30-day: 48 origins each
  split (each forecast must fit inside its evaluation window).
\item
  Four candidates per horizon: seasonal naive, additive weekly ETS,
  XGBoost autoregression (28-day lags), MI/PCA-selected Ridge ---
  selected independently per horizon by validation MAE only.
\end{itemize}
\end{block}
\end{frame}

\begin{frame}[fragile]
\begin{block}{7. Validation selection \& final test}
\protect\phantomsection\label{7-validation-selection--final-test}
\begin{itemize}
\tightlist
\item
  \textbf{7-day:} XGBoost wins validation (MAE 4.181) over ETS, Ridge,
  naive. Test MAE \textbf{5.05} vs seasonal baseline 6.41.
\item
  \textbf{30-day:} weekly-additive ETS wins validation (MAE 4.394). Test
  MAE \textbf{4.67}, RMSE 7.20, WAPE 27.5\%, sMAPE 41.3\%, 90\% interval
  coverage 87.7\%, vs seasonal baseline test MAE 6.49.
\item
  Selection is \textbf{never} changed on test rank ---
  ETS\textquotesingle s lower 7-day test MAE does not override the 7-day
  XGBoost validation selection.
\item
  Bounded hyperparameter sensitivity/tuning search re-validates any
  candidate on the honest rolling-origin backtest; current baseline
  confirmed as the best searched configuration
  (\texttt{keep\_baseline}).
\end{itemize}
\end{block}
\end{frame}

\begin{frame}
\begin{block}{8. Metrics, uncertainty \& explainability}
\protect\phantomsection\label{8-metrics-uncertainty--explainability}
\begin{itemize}
\tightlist
\item
  MAE/RMSE describe absolute daily miss (RMSE penalizes larger misses);
  WAPE/sMAPE express relative error. 90\% intervals are
  validation-residual quantiles --- measured test coverage
  (86.6\%/87.7\%) runs below nominal 90\%, so calibration is imperfect
  and not guaranteed to persist.
\item
  Challenger-model explanations (pre-validation XGBoost fit only):
  \textbf{SHAP} (leading lags 28/21/14), \textbf{PDP} (average marginal
  effect, top 3 features), \textbf{ICE} (per-validation-day curves
  exposing heterogeneity PDP averages mask), \textbf{LIME} (one
  representative local explanation).
\item
  Primary 30-day ETS is interpreted separately through its own
  level/trend/ weekly-seasonal components and residual diagnostics ---
  none of the XGBoost explanations above explain the ETS artifact.
\end{itemize}
\end{block}
\end{frame}

\begin{frame}
\begin{block}{9. Ethics \& fairness interpretation}
\protect\phantomsection\label{9-ethics--fairness-interpretation}
\begin{itemize}
\tightlist
\item
  Source has no approved protected demographic attributes (age, gender,
  ethnicity, etc.) --- \textbf{protected-group fairness cannot be
  assessed}, not because it was skipped, but because the data to run it
  does not exist.
\item
  Substituted operational-equity proxy:
  priority/year/masked-project-group breakdowns with \textbf{Wilson
  confidence intervals} for small-group rates.
\item
  Standard mitigation toolkit (reweighting, thresholds, augmentation,
  post-processing) is classification/protected-attribute tooling that
  does not transplant onto a no-protected-attribute
  regression/forecasting task; substituted safeguards are the
  operational audit above, a recommendation to collect governed
  protected-attribute data before any fairness claim, and mandatory
  human review of forecast-driven decisions.
\end{itemize}
\end{block}
\end{frame}

\begin{frame}[fragile]
\begin{block}{10. Integration readiness}
\protect\phantomsection\label{10-integration-readiness}
\begin{itemize}
\tightlist
\item
  \textbf{Locally validated:} all eight notebooks executed; 52/52 tests
  pass without skips; FastAPI demo passed a live localhost
  health/forecast/ validation check; MLflow run logs parameters and
  metrics.
\item
  \textbf{Serving scope:} the API serves one persisted primary 30-day
  ETS artifact for all accepted horizons --- \texttt{horizon\_days=7}
  does \textbf{not} load the validation-selected 7-day XGBoost model.
\item
  \textbf{Not ready for operational use:} unresolved date-range
  provenance mismatch; unverified zero-arrival completeness assumption;
  dataset ends March 2023; SLA mapping/clock unapproved; no current
  feed; no real-world staffing outcomes or costs measured.
\end{itemize}
\end{block}
\end{frame}

\begin{frame}
\begin{block}{11. Conclusion \& next steps}
\protect\phantomsection\label{11-conclusion--next-steps}
\begin{itemize}
\tightlist
\item
  Validation selects XGBoost at 7 days and ETS at 30 days; both beat the
  seasonal-naive baseline on held-out test MAE. Neither this result nor
  lower historical error establishes operational or business impact.
\item
  \textbf{Before operational reliance:} resolve the provenance/coverage
  caveats; obtain current representative arrivals beyond March 2023;
  confirm SLA clock/policy with service owners; agree business
  tolerances in advance; test on a genuinely future period; run a
  human-reviewed ROI pilot with approved cost/outcome data.
\item
  \textbf{Bottom line:} a corrected, reproducible historical forecasting
  prototype with a disclosed retrospective SLA baseline --- not a
  contractual SLA-compliance system or a proven staffing-cost/ROI tool.
\end{itemize}
\end{block}
\end{frame}

\begin{frame}[fragile]
\begin{block}{Format note}
\protect\phantomsection\label{format-note}
This file (\texttt{presentations/technical\_slides.md}) renders two ways
via \texttt{src/generate\_technical\_reveal\_deck.py}:

\begin{itemize}
\tightlist
\item
  \textbf{reveal.js HTML}
  (\texttt{presentations/technical\_reveal\_deck.html}) --- open
  directly in any browser; no install required beyond an internet
  connection to load the reveal.js CDN assets referenced by the page.
\item
  \textbf{Beamer LaTeX source}
  (\texttt{presentations/technical\_reveal\_deck\_beamer.tex}) --- a
  valid Beamer \texttt{.tex} source file. No LaTeX engine is installed
  in this environment, so it is provided as compileable source, not a
  pre-rendered PDF; compile it with \texttt{pdflatex}/\texttt{xelatex}
  where available.
\end{itemize}

\texttt{technical\_capstone.pptx} and
\texttt{technical\_presentation.pptx} in this folder are unchanged by
this export.
\end{block}
\end{frame}

\end{document}
